\begin{prop}\label{prop:no-cyc}
A triangulation $\mathcal{T}$ contains an interior $(d+1)$-simplex if and only if $Q(\mathcal{T})$ contains a cycle.
\end{prop}
\begin{proof}
We first suppose for contradiction that $\mathcal{T}$ contains no interior $(d+1)$-simplices and that $Q(\mathcal{T})$ does contain a cycle. By Lemma~\ref{lem:single-step-arrows}, we can realise this cycle as a path $P \colon A \leadsto B$ in $\tqdn$, where $A$ is some vertex in the cycle in $Q(\mathcal{T})$ and $\pi(B) = A$ with $A \neq B$. We consider this path $P$ as a subquiver of $\tqdn$. By Lemma~\ref{lem:all_paths}, every path $A \leadsto B$ in $\tqdn$ must lie in $\overline{P}$. There is a path $A \leadsto A + \onevec \leadsto B$ in $\tqdn$. Hence $A + \onevec$ is a vertex of $\overline{P}$. By Lemma~\ref{lem:simp_arrow_switch}, if $C$ is a vertex of $\overline{P}$, then $\pi(C)$ is a vertex of $Q(\mathcal{T})$. Therefore $\pi(A + \onevec)$ is a vertex of $Q(\mathcal{T})$, but this is a contradiction, since $\pi(A + \onevec)$ and $A$ are intertwining.

We now suppose that $\mathcal{T}$ is a triangulation with an interior $(d + 1)$-simplex  $(a_{0}, \dots, a_{d + 1})$. Then $Q(\mathcal{T)}$ has a cycle given by concatenating the paths
\begin{align*}
(a_{0}, a_{1}, \dots, a_{d - 1}, a_{d}) &\leadsto (a_{0}, a_{1}, \dots , a_{d - 1}, a_{d + 1}) \leadsto (a_{0}, a_{1}, \dots, a_{d - 2}, a_{d}, a_{d + 1}) \leadsto \\
 \dots &\leadsto (a_{1}, a_{2}, \dots, a_{d}, a_{d + 1}) \leadsto (a_{0}, a_{1}, \dots , a_{d}),
\end{align*}
noting Lemma~\ref{lem:find_path}.
\end{proof}
